\documentclass[11pt, twocolumn]{article}
\usepackage[super,comma,sort&compress]{natbib}

\begin{document}


%%%%%%%%
%Part 1: Literature Review
%%%%%%%%

\part{Literature Review}
	\section{Introduction}
		When we talk about the concept of ``Thinking", there're many theory that try to elaborate in which method is the best. In this Week we'll only be focusing on 2 types of method:
		\begin{enumerate}
			\item Linear Thinking
			\item Design Thinking
		\end{enumerate}
	
	\section{Linear Thinking}
		\subsection{What is Linear Thinking?}
			Linear Thinking is often related to scientist perspective on thinking. The core idea of Linear Thinking is as the name suggested ``Linear".
			This method of thinking is a way of reasoning that follow a step\textendash by\textendash step path from a problem to solution. The example is as shown below:
				\begin{enumerate}
					\item Traffic Congestion in Bangkok
						\begin{itemize}
							\item Build more roads, and widen highways.
						\end{itemize}
					\item Traffic will flow smoothly.
				\end{enumerate}
			As you can see, the path in which solutions are form is in linear path: Problem to Solution.
	
	\section{Design Thinking}
		\subsection{What is Design Thinking?}
			Design Thinking is a way of approaching problems that focus on creativity; and accepts that problems are often unclear and complex which is refer as ``wicked problem".
		\subsection{Wicked Problem PT. I}
		To understand what ``Wicked Problem" are, we have to first understand the nature of it.
			\subsubsection{10 Characteristics of Wicked Problem.}
				\begin{enumerate}
					\item Wicked problems have \textbf{no definitive formulation}, but every formulation of a wicked problem corresponds to the formulation of a solution.
					\item Wicked problems have \textbf{no stopping rules}.
					\item Solutions to wicked problems \textbf{cannot be true or false, only good or bad}.
					\item  In solving wicked problems there is \textbf{no exhaustive list of admissible operations}.
					\item For every wicked problem \textbf{there is always more than one possible explanation, with explanations depending on the \textit{Weltanschauung}\footnote{Weltanschauung identifies the intellectual perspective of the designer as an integral part of the design process.} of the designer.}
					\item Every wicked problem is \textbf{a symptom of another, ``higher level" problem.} 
					\item No formulation and solution of a wicked problem has a definitive test.
					\item Solving a wicked problem is a ``one shot" operation, with no room for trial and error.
					\item Every wicked problem is unique.
					\item The wicked problem solver has no right to be wrong-they are fully responsible for their actions
				\end{enumerate}
\newpage
			\subsubsection{Interpreting the 10 Characteristics.}
				\begin{enumerate}
					\item There are no exact way to solve the wicked problem.
					\item There are no clear end point of wicked problem.
					\item When solving a wicked problem; there are no correct, or non-correct solution. There are only good or bad solution.
					\item There are no fix method to fix wicked problem.
					\item Each of the explanation (for wicked problem) will be based on how the designer perceive those problem.
					\item Wicked problems always emerge from another, and bigger, problem.
					\item Solutions of the wicked problem cannot be test.
					\item You can only solve wicked problem once per problem.
					\item Every wicked problem will not be the same.
					\item By solving the wicked problem, the solver, or ``designer" will have to bear the responsibility of the outcome.
				\end{enumerate}
				
		\subsection{Wicked Problem PT. II}
		
			The 10 Characteristics of Wicked Problem state clearly on what nature of this type of problem is. However, it still left an unanswered question:
			\begin{center}
				 \textbf{WHY ARE DESIGN PROBLEMS INDETERMINATE AND, THEREFORE, WICKED?}
			\end{center}
			To easily explain why design problems are wicked, we should explore the definition of it:
			\textbf{A WICKED PROBLEM IS PROBLEMS THAT CAN BE SOLVED IN VARIOUS FIELD OF STUDY, AND THE SOLUTIONS TO WICKED PROBLEM WILL ADAPT AND CHANGE ACCORDING TO WHAT THE SOLVER PERCEIVES THE PROBLEM TO BE.}
%%%%%%%%
%Part 2: Richard Buchanan on Design Thinking
%%%%%%%%
\part{Richard Buchanan on Design Thinking}
	In the article by \textit{Richard Buchanan}, the author aims to raise the important of using design thinking in the modern world.
	
	As review in the \textbf{Literature Review} section. The usage of Linear Thinking is not enough to truly solve the problem in society, and by understanding how Design Thinking work, and utilise it; we'll be able to move forward to the better future.
	
	\textbf{TLDR; Design Thinking method should replace Linear Thinking, and be the modern way of thinking.}
	
\end{document}
